% =============================================================================
\section{Related Work}
\label{sec:related}
% =============================================================================

\paragraph{Membership inference attacks.}
Shokri et al.~\cite{shokri2017membership} introduced the shadow-model
attack and established Purchase100 and Texas100 as standard benchmarks.
Subsequent work strengthened attacks via label-only
access~\cite{choquette2021label}, enhanced
thresholds~\cite{ye2022enhanced}, and the likelihood-ratio
formulation~\cite{carlini2022lira} that we use as ground truth.
Salem et al.~\cite{salem2019ml} showed that attacks are robust to
imperfect shadow model knowledge.
Our work does not propose a new attack. Rather, it ranks datasets by their
attack success from data properties alone, before any model is trained.

\paragraph{Privacy risk theory.}
Yeom et al.~\cite{yeom2018privacy} proved that MIA advantage is bounded
by the generalization gap, linking privacy to overfitting.
Sablayrolles et al.~\cite{sablayrolles2019whitebox} characterized the
Bayes-optimal MIA as a likelihood ratio test, which forms the basis of
our Theorem~\ref{thm:main}.
Our contribution extends these results by showing that the likelihood
ratio, and hence the MIA advantage, is bounded, under explicit modeling
assumptions, by pre-training distributional statistics, not only by
post-training model properties.

\paragraph{Memorization and long-tail data.}
Feldman~\cite{feldman2020memorization} proved that learning algorithms
must memorize atypical (long-tail) training examples to achieve near-optimal
generalization, and connected this to per-sample privacy risk.
Carlini et al.~\cite{carlini2021extracting} showed that LLMs memorize
rare training examples verbatim.
Our work differs along three dimensions.
Unit of analysis: we study dataset-level statistics rather than
per-sample certificates.
Timing: we operate before training, with no trained model.
Output: a dataset-level risk score used to rank datasets, not a post-hoc
per-sample memorization label.
Our results are consistent with Feldman's. The datasets we find most at
risk are those with high sample uniqueness, the aggregate counterpart of
atypical, memorization-prone examples. We complement that line of work by
asking, at the dataset level and before training, how far this geometry
predicts risk, and we find that it does so moderately but significantly.

\paragraph{Differential privacy.}
DP-SGD~\cite{abadi2016deep} is the standard mechanism for training
private neural networks.
Jagielski et al.~\cite{jagielski2020auditing} show empirically that the
privacy guarantees of DP-SGD are tight.
Zanella-Béguelin et al.~\cite{zanella2023bayesian} propose Bayesian
estimation of the actual $\varepsilon$ achieved.
None of these works ask how the benefit of a given $\varepsilon$
varies with dataset structure. We raise this as a direction motivated by our
bound (Section~\ref{sec:implications}), though we do not evaluate DP-SGD in
this work.

\paragraph{Dataset auditing and data valuation.}
Long et al.~\cite{long2020pragmatic} propose a pragmatic approach to
membership inference that accounts for dataset characteristics.
Our work goes further by predicting risk before training and providing a
formal theoretical basis through Theorem~\ref{thm:main}.
